\section*{第 \ref{ch:b2:scalar-light-model} 章——光的标量模型}

\begin{solution}{exo:b2:scalar-light-model:1}
$(AB) = \qty{1.5}{cm}$；$25\,000$ 个波长；比空气中多 \qty{0.5}{cm}，即 $8300$
个波长，$\Delta\varphi = 2\pi \times 8300$，延迟 \qty{17}{ps}；产生一个波长
延迟所需的厚度 $e = \lambda/(n - 1) = \qty{1.2}{\micro m}$。
\end{solution}

\begin{solution}{exo:b2:scalar-light-model:2}
$\ell_c = \lambda^2/\Delta\lambda$ 或 $c/\Delta\nu$，$\tau_c = \ell_c/c$：白光 \qty{0.9}{\micro m}（\qty{3}{fs}）；
钠灯 \qty{1.7}{cm}（\qty{60}{ps}）；发光二极管 \qty{16}{\micro m}（\qty{50}{fs}）；二极管 \qty{3}{mm}
（\qty{10}{ps}）；稳频激光器 \qty{300}{km}（\qty{1}{ms}）。\qty{1}{mm}：钠灯、
二极管和激光器；\qty{10}{m}：仅稳频激光器。
\end{solution}

\begin{solution}{exo:b2:scalar-light-model:3}
$10^{-3}/3.06 \times 10^{-19} = \qty{3.3e15}{}$ 每秒。灯泡：$I = 5/4\pi \times 9 =
\qty{44}{mW/m^2}$，瞳孔 \qty{28}{mm^2}：\qty{1.2}{\micro W}，$3.5 \times 10^{12}$ 个光子每秒，
$3.5 \times 10^{11}$ 个（在 \qty{0.1}{s} 内），涨落为 $2 \times 10^{-6}$：完全平滑。
\end{solution}

\begin{solution}{exo:b2:scalar-light-model:4}
精确值 $\sqrt{1 + 10^{-4}} - 1 = \qty{49.9988}{\micro m}$，近轴值 \qty{50}{\micro m}：约合
一百个波长；波前是半径为 \qty{1}{m} 的球面；
四次项 $r^4/8d^3$ 达到 $\lambda$ 时，$r = (8\lambda d^3)^{1/4} = \qty{4.5}{cm}$。
\end{solution}

\begin{solution}{exo:b2:scalar-light-model:5}
(a) 半径为 $f$、以 $F'$ 为球心的球面，在透镜平面上
相对光轴的相位为 $-\pi r^2/\lambda f$（近轴近似），恰好是
透镜的相位。(b) 总相位为 $\pi r^2/\lambda\,(1/d - 1/f) = -\pi r^2/\lambda d'$：一列波
会聚于 $d'$，且 $1/d' = 1/f - 1/d$。(c) 玻璃引入的相位为 $(n - 1)e(r)
\times2\pi/\lambda$：$e(r) = e_0 - r^2/2(n - 1)f$；矢高 $10^{-4}/0.1 = \qty{1}{mm}$。(d) 按模 $2\pi$ 定义的
相位仍是同一相位——在设计
波长下如此；在其他波长下，台阶不再是整数个波长。
\end{solution}

\begin{solution}{exo:b2:scalar-light-model:6}
(a) 自 $x = -D$ 处的平面波前出发，经 $(x, y)$ 再到 $F(f, 0)$：$D + x
+ \sqrt{(x - f)^2 + y^2} = D + f$，于是 $\sqrt{(x - f)^2 + y^2} = f - x$，$y^2 = 4fx$，
其中 $F$ 为焦点。(b) 圆 $x = 2f - \sqrt{4f^2 - y^2} \approx y^2/4f +
y^4/64f^3$ 仅在二阶上相符。(c) $y^4/64f^3 = \qty{1.6}{\micro m}$，经反射加倍为
\qty{3}{\micro m}，相当于六个波长，因此需采用抛物
面镜。(d) 在 $\lambda \sim \qty{1}{cm}$ 时，同样的误差仅为
一个波长的千分之一。
\end{solution}

\begin{solution}{exo:b2:scalar-light-model:7}
(a) $L(x) = n_1\sqrt{a^2 + x^2} + n_2\sqrt{b^2 + (d - x)^2}$：$L' = n_1\sin\theta_1 -
n_2\sin\theta_2 = 0$。(b) 在驻定路径附近，相邻路径的
相位相等，其波叠加；在别处则相互抵消。(c) 若 $n_1 = n_2$，
则 $\sin\theta_1 = \sin\theta_r$。(d) 严格成像意味着从 $A$ 到 $A'$ 的所有光线
光程相等：驻定路径是一整族，而不是一条。
\end{solution}

\begin{solution}{exo:b2:scalar-light-model:8}
(a) $I = 2I_0(1 + \cos\Delta\varphi)$：$4I_0$ 与 $0$。(b) $N = 10^7$：平均为 $2I_0$，残余
约为其 $\sim 2I_0/\sqrt N = 6 \times 10^{-4}$ 倍。(c) 在 \qty{10}{ps} 内相位
保持不变：探测器会看到条纹，每 \qty{0.1}{ns} 跳变一次。(d) 对于
$\tau_c \sim \qty{1}{\micro s}$ 的激光器，若探测器比之更快，则条纹
存在一微秒，并已被拍摄。
\end{solution}

\begin{solution}{exo:b2:scalar-light-model:9}
(a) 小角近似：$\tan i \approx n\tan r$，光线似乎来自 $h/n$ 处。
(b) \qty{0.90}{m}；光程为 $n \times 1.2 = \qty{1.6}{m}$，多出 \qty{0.4}{m}。
(c) $n \times 1 = \qty{1.33}{m}$。(d) 在大角度下折射更强，
视深更小。
\end{solution}

\begin{solution}{exo:b2:scalar-light-model:10}
(a) 两列等幅波，相位差为 $2\pi\nu\delta/c$。(b) 在
整个频带上取平均：
\[
\frac1{\Delta\nu}\int\cos\frac{2\pi\nu\delta}c\,\dd\nu = \cos\frac{2\pi\nu_0\delta}c\,\operatorname{sinc}\frac{\pi\Delta\nu\,\delta}c .
\]
(c)
$I_{\max, \min} = 2I_0(1 \pm |\operatorname{sinc}|)$：$V = |\operatorname{sinc}(\pi\Delta\nu\delta/c)|$，
在 $\delta = c/\Delta\nu$ 处为零。(d) $\Delta\nu = c\Delta\lambda/\lambda^2 = \qty{1.9e13}{Hz}$：$\delta =
\qty{16}{\micro m}$。
\end{solution}

\begin{solution}{exo:b2:scalar-light-model:11}
(a) 在 \qty{38}{mm^2} 上每秒 $100$ 个光子：\qty{2.6e6}{m^{-2}.s^{-1}}，$\times4 \times
10^{-19}$ J：\qty{1e-12}{W/m^2}。(b) $10^{-13.1} \times 1000 = \qty{8e-11}{W/m^2}$，$2 \times 10^8$
个光子（每平方米每秒），每秒进入瞳孔的为 $8000$ 个
——远高于 (a)：眼睛的实际极限由背景与
众多视杆细胞上的分布共同决定。(c) $(200/7)^2 = 800$，$+7.3$ 个星等：$m \approx
13$。(d) $N$ 个光子涨落为 $\sqrt N$：信噪比随
曝光量的平方根增长。
\end{solution}

\begin{solution}{exo:b2:scalar-light-model:12}
(a) $d(1 + r^2/2d^2 - r^4/8d^4 + \dots) - d$。(b) $r^4/8d^3 < \lambda/4$。(c) $r <
(2\lambda d^3)^{1/4}$：\qty{3.2}{cm}（对 \qty{1}{m}）；\qty{1}{mm}（对 \qty{1}{cm}），角度为
\qty{0.1}{rad}。(d) 球差即此四次项；其余
像差都是它的离轴“同类”。
\end{solution}

\begin{solution}{pb:b2:scalar-light-model:1}
\textbf{1.} $\nu_1 = \qty{5.0934e14}{Hz}$，$\nu_2 = \qty{5.0882e14}{Hz}$：$\Delta\nu = \qty{5.2e11}{Hz}$。

\textbf{2.} $c/\Delta\nu_{\text{自然}} = \qty{30}{m}$。

\textbf{3.} $\sqrt{k_BT/m} = \qty{425}{m/s}$：$\Delta\nu_D \approx 2\nu v/c \approx \qty{1.4}{GHz}$，约为
自然线宽的一百倍。

\textbf{4.} $c/\Delta\nu_D \approx \qty{20}{cm}$；计入碰撞后，约为 \qty{2}{cm}。

\textbf{5.} $\lambda^2/\Delta\lambda = 589^2/0.6$ nm $= \qty{0.58}{mm}$：从重合
到反相每隔 \qty{0.29}{mm} 一次；每拍约 $1000$ 条条纹。

\textbf{6.} 对比度在光程差上以 \qty{0.58}{mm} 为周期振荡
（零点位于 $0.29$、$0.87$、\qty{1.45}{mm}、\dots），并受一个包络线
约束：该包络在几厘米内逐渐衰减，约五十次拍之后便什么也没有。

\textbf{7.} 每秒 \qty{3e18}{} 个光子；\qty{2.1}{eV}，属于
光谱中的黄色。

\textbf{8.} 不能：原子彼此独立，相位随机。同一灯的
两条谱线同样来自不同原子，彼此之间不发生干涉。
拍即两组条纹系统在\emph{强度}上叠加。

\textbf{9.} $\lambda^2/\Delta\lambda = \qty{0.85}{mm}$；$c/\Delta\nu = \qty{300}{m}$。

\textbf{10.} $\qty{3.3e15}{}$ 每秒，沿光束以 $c/N = \qty{90}{nm}$ 为间距排布：
相干长度内有一万个光子。干涉并非光子相遇，
而是每个光子与自身发生干涉。

\textbf{11.} $\Delta\nu = c\Delta\lambda/\lambda^2 = \qty{360}{GHz}$；纵模间距 $c/2nL = \qty{43}{GHz}$：
约八个纵模。

\textbf{12.} $e \times 3.3 \times 10^{15} = \qty{0.5}{mA}$。

\textbf{13.} \qty{2}{mm} $>$ \qty{0.85}{mm}：无条纹；\qty{10}{m} $<$ \qty{300}{m}：有条纹。

\textbf{14.} 一个纳秒脉冲长 \qty{30}{cm}，远长于
\qty{0.85}{mm} 的相干长度：决定相干性的是谱线宽度，而非脉冲本身，
相继脉冲只有在关断期间激光器的相位仍然保持时才会发生干涉，
而它并不能。

\textbf{15.} 粗糙墙面散射的相干光在视网膜上
因随机光程差而相互干涉：形成散斑。灯的
微米量级的相干性把这一切统统平均掉。

\textbf{16.} $\lambda^2/\Delta\lambda = \qty{17}{\micro m}$：远不足以在粗糙墙面上
形成散斑——因此没有散斑，图像更平滑。

\textbf{17.} $(n - 1)e = \qty{1}{mm}$：$1540$ 条干涉条纹。

\textbf{18.} $r = \qty{6.65}{\degree}$；附加光程 $\approx e(n/\cos r - \cos(i - r)/\cos r)
= 3.021 - 2.010 = \qty{1.011}{mm}$，比未倾斜时多 \qty{11}{\micro m}：约多出
二十条干涉条纹。

\textbf{19.} 每条条纹占五个像素；若每条条纹仅对应一个像素，则该像素
将积分整整一个周期，对比度随之消失。

\textbf{20.} 像素上的功率 \qty{1e-12}{W} $\times$ \qty{10}{ms} $= \qty{1e-14}{J}$：$3 \times 10^4$
个光子，噪声为 $0.6\%$。

\textbf{21.} $I = 2I_0[1 + \cos(2\pi\delta/\lambda)]$：光程差每变化一个波长，
出现一条条纹。

\textbf{22.} 人眼把千赫兹的运动平均为
一个均匀场；光电二极管则跟随这一运动。

\textbf{23.} $(n - 1) \times 10 = \qty{2.9}{mm}$：$4500$ 条干涉条纹；压强变化 $1\%$
会使条纹移动 $45$ 条——除非抽真空，否则
干涉仪就是一台气压计。

\textbf{24.} 光子数率几乎相同（\qty{2.1}{eV} 与 \qty{1.9}{eV} 相比）；
条纹取决于振幅与相位，而光子计数仅
与噪声有关。

\textbf{25.} 灯：厘米量级，受多普勒与碰撞展宽限制，
可见条纹达约 \qty{1}{cm}；廉价激光笔：\qty{0.85}{mm}，受其多个纵模限制；
稳频激光器：\qty{300}{m}，受其
谱线宽度限制。
\end{solution}
